\documentclass{standalone}
%\immediate\write18{mkdir -p figures}

\usepackage{tikz}
\usepackage{xcolor}
\usetikzlibrary{external}
\usetikzlibrary{fadings}
\tikzexternalize[prefix=cabecalho/] % Cria PDFs externos na pasta 

\begin{document}
\begin{tikzpicture}
%%%%%%%%%%%%%%%%%%%%%%%%%%%%%%%%%%%%%
%%%% DEFINA A COR BASE DA FIGURA %%%%
\definecolor{cormodelo}{RGB}{163, 11, 11}
%%%%%%%%%%%%%%%%%%%%%%%%%%%%%%%%%%%%%


% Parâmetros
\def\rows{7}
\def\cols{15}
\def\size{2}

% Loop para desenhar triângulos e quadrados com tons de azul
\foreach \i in {0,...,\cols} {
  \foreach \j in {1,...,\rows} {
    % Define a cor com base na altura: do azul ao branco
    \pgfmathsetmacro{\fade}{\j/\rows*100}

    \ifodd\numexpr\j\relax
    %Ponto inferior esquerdo
    \pgfmathsetmacro{\x}{\i*\size}
    \pgfmathsetmacro{\y}{(\j-1)*\size}
    
    \pgfmathsetmacro{\basePct}{(\j/\rows)*80+5}
    \pgfmathsetmacro{\shadePct}{max(min(\basePct+(rnd*10-20),100),0)+7*rnd}
      
      \filldraw[cormodelo!\shadePct, cormodelo!\shadePct] (\x,\y) -- ({\x+\size},\y) -- ({\x+\size},{\y+\size}) -- cycle;

    \pgfmathsetmacro{\basePct}{(\j+0.5)/\rows*80+5}
    \pgfmathsetmacro{\shadePct}{max(min(\basePct+(rnd*10-20),100),0)+7*rnd}
      
      \filldraw[cormodelo!\shadePct, cormodelo!\shadePct] (\x,\y) -- (\x,{\y+\size}) -- ({\x+\size},{\y+\size}) -- cycle;
      
      \else
    %Ponto superior esquerdo
      \pgfmathsetmacro{\x}{\i*\size}
    \pgfmathsetmacro{\y}{(\j*\size}
    
    \pgfmathsetmacro{\basePct}{(\j+0.5)/\rows*80+5}
      \pgfmathsetmacro{\shadePct}{max(min(\basePct+(rnd*10-20),100),0)+7*rnd}
      
      \filldraw[cormodelo!\shadePct, cormodelo!\shadePct] (\x,\y) -- ({\x+\size},\y) -- ({\x+\size},{\y-\size}) -- cycle;

    \pgfmathsetmacro{\basePct}{(\j)/\rows*80+5}
    \pgfmathsetmacro{\shadePct}{max(min(\basePct+(rnd*10-20),100),0)+7*rnd}
      
      \filldraw[cormodelo!\shadePct, cormodelo!\shadePct] (\x,\y) -- (\x,{\y-\size}) -- ({\x+\size},{\y-\size}) -- cycle;
      \fi
  }
}

\fill[color=white, path fading=north] (0,-1) -- (\cols*\size+\size,-1) -- (\cols*\size+\size,10) -- (0,10);
\end{tikzpicture}
\end{document}
